\section{A height arch at a constant boundary arc}\label{sec:boundary}

We now recover the boundary information lost in the single-color windows.
Near an arc with prescribed height $k$, we construct a path of height at
least $k$ that joins the boundary on both sides of a small exceptional
piece.  Closing this path through the exterior separates that piece from
the rest of the domain.  The estimate uses the actual conditional Gibbs
law and will also apply to the new boundary traces exposed during this
separation.  In \cref{sec:limits}, lower and upper versions identify both
the boundary mean and its mixed moments with distant observations.

\paragraph{Regions, fixed sites, and artificial exterior.}
All graphs in this section are embedded in the triangular lattice of mesh
$\delta$.  A \emph{skeleton} is a finite connected union of lattice edges.
The free sites of a region lie in bounded faces of its skeleton.  Every
neighbor of a free site is either free or a fixed skeleton site; the latter
sites are its \emph{usable boundary}.  Conditional on the values $\eta$ at
the usable boundary, the height is uniform among the admissible integer
Lipschitz extensions.  We always assume that this set is nonempty.
Only usable boundary values enter this specification.

Initially the skeleton $S_0$ is a polygon boundary.  We also allow it to be
the union of that boundary and finitely many simple lattice circuits meeting
it.  The additional circuits cut off small boundary patches; the chosen
remaining free regions are on the other side.  Vertices in the removed
patches, and vertices exterior to the original polygon, will be declared
open artificially.  Usable boundary sites are never declared open
artificially.  Thus, at threshold $a$, the open sites are
\begin{equation}\label{eq:boundary-open-sites}
 \{v\text{ free or usable}:H(v)\ge a\}
 \ \cup\ \{v\text{ outside the region and its usable boundary}\}.
\end{equation}
When an additional separator is exposed, its usable sites carry their
actual exposed heights.  In particular, this convention does not discard a
new boundary condition.

Write $Q_r=p+[-r,r]^2$ in tangent and inward-normal coordinates at a smooth
boundary point $p$.  The interior is on the positive normal side.
An open site circuit surrounding $Q_r$ is called an \emph{arch separator}:
its part in the domain joins the clean boundary on the two sides, and its
remaining part can run through the artificial exterior.  It separates the
free sites near $p$ from remote boundary data.

\begin{theorem}[Constant-arc arch estimate]\label{thm:boundary-arches}
There are constants $M>1$, $c,C>0$, and a positive geometric tolerance with
the following property.  Fix a positive integer $N$ and scales
\begin{equation}\label{eq:arch-scales}
  16r<\ell_1<\cdots<\ell_N<R/10,
  \qquad \ell_{i+1}\ge M\ell_i.
\end{equation}
Suppose that, in $Q_R$, the limiting boundary is the single graph
$y=f(x)$, with $f(0)=f'(0)=0$ and $\|f'\|_\infty$ smaller than the
geometric tolerance, and that its inward side is above this graph.  The
original polygon boundary is a sufficiently fine uniform parametric
approximation of this arc and has the same inside and outside away from
its shrinking boundary neighborhood.  Its approximation error, divided by
$\ell_1$, tends to zero.  A smooth boundary has these properties after
decreasing the fixed outer scale.
Allow, in addition, a fixed finite number of boundary-attached patch
replacements as above, whose diameters are sufficiently small relative to
the selected scales.  Suppose that the actual usable boundary data satisfy
\begin{equation}\label{eq:arch-boundary-data}
 \eta\ge k-2\quad\hbox{everywhere},
 \qquad
 \eta\ge k\quad\hbox{on the usable boundary in }Q_R\setminus Q_r,
\end{equation}
where $k\in\Z$.  For sufficiently small mesh and sufficiently fine boundary
and replacement approximations,
\begin{equation}\label{eq:arch-high-probability}
 \PP_\eta\bigl(\text{an arch at threshold }k\text{ surrounds }Q_r
       \text{ in }Q_R\bigr)\ge 1-Ce^{-cN}.
\end{equation}
The constants do not depend on the boundary heights or on the number of
vertices in a boundary or separator.  They also do not depend on the fixed
number of tiny replacements; the required approximation accuracy may do so.
There is a reversed statement with $\eta\le k+2$, clean data $\eta\le k$,
and a separator at height at most $k$.
\end{theorem}

The order of quantifiers is part of the statement.  One first chooses the
number and the positive scales of the comparison boxes.  The mesh and the
parametric boundary error then tend to zero, and the diameters of the fixed
finite family of replacements tend to zero relative to these boxes.
Consequently, the error tends to zero as the number of available scales
between a central exception and the end of a clean arc tends to infinity.
There is no assertion uniform over unresolved geometry at the mesh scale.
The lower estimate has no upper bound on $\eta$ as a hypothesis.  It
therefore applies to bounded boundary data and also, conditionally, to the
possibly larger random values revealed on a preliminary separator.

The proof first constructs many pairs of weaker separators, at height
$k-1$.  Selecting the inner separator from its inner side and the outer
separator from its outer side leaves an unexamined Gibbs region between
them.  These retained middle regions are conditionally independent.  In
each one we then build two attachments from the clean boundary at the
required height $k$, and join them by bulk paths.  A fixed positive chance
of this upgrade in each retained region gives the exponential error in
\eqref{eq:arch-high-probability}.

The attachment is the main step: a synthetic mean comparison first detects
the high side of a boundary-entering level contour; two low lateral
barriers then force the adjacent path at height $k$ to rise into the domain.
Closed high islands alone would not give an attachment.  All boxes in this construction
have fixed positive relative margins.  After they have been chosen, the
mesh is small enough that rounding preserves their separation and
intersection properties.

\subsection{An anchored two-interval comparison}

For any $k$, one of the two Gray-code colors, denoted by $t$, and one of its
signs $s$, describe the two adjacent residues in $[k-1,k]$.  The opposite
sign of the same color describes $[k-3,k-2]$.  Denote the complementary
color by $v$.  The transition from the low interval to the high interval
can only be the step
\begin{equation}\label{eq:anchored-transition}
                  k-2\longleftrightarrow k-1.
\end{equation}
The complementary color has the same value, say $c_k$, at both ends of
this step, and this is the same value at every such transition.

Here is the comparison we shall repeatedly use.  On a connected skeleton
$S$, assign the sign $s$ to selected connected portions and $-s$ elsewhere.
Require $v=c_k$ at each of a bounded number $b$ of transition edges, and
anchor the lift by \eqref{eq:anchored-transition} at one of them.  A
constant-$t$ chain stays in one adjacent pair of integers: a nearest-neighbor
step cannot move to the next pair with the same residue color.  A transition
with the prescribed complementary color fixes the next pair.  Connectedness
therefore gives the exact bounds
\begin{equation}\label{eq:synthetic-intervals}
 K\in[k-1,k]\text{ on the high portions},\qquad
 K\in[k-3,k-2]\text{ on the low portions}.
\end{equation}
Extra lattice edges between pinned vertices only impose additional
compatibility; they do not introduce new branches of the lift.  In
particular the prescription is nonempty: taking $v\equiv c_k$ gives just
the two values in \eqref{eq:anchored-transition}.

Construct this synthetic height in a larger, finite, simply connected free
lattice volume containing $S$ and all comparison boxes.  Given $t$, the
complementary color is constant on each component of $t$-changing edges and
its component labels are independent fair signs.  The transition
requirements have conditional probability at least $2^{-b}$.  The same
lower bound holds if some transition sites belong to one component, since
all requested labels are $c_k$.

Suppose the actual boundary is at least $k$ on the selected high portions
and at least $k-2$ elsewhere.  Every possible synthetic usable trace is then
bounded above by the actual trace.  Conditional on its trace, the synthetic
interior is a uniform height Gibbs law.  Height comparison, which follows
from the min/max lattice condition, gives
\begin{equation}\label{eq:synthetic-height-domination}
       K|_{\text{free sites}}\preceq H|_{\text{free sites}}.
\end{equation}
This compares the interiors with their actual fixed traces; no observation
of an unconditioned boundary is substituted for boundary pinning.

The prescriptions used below have at most four transitions on the base
polygon.  For tiny replacements, choose those transition points away from
the attachment neighborhoods and give an entire replacement the local
prescription of its base arc.  Its branch is then anchored through that arc
and it adds no transition.  To justify this uniformly, fix the finitely many
selected boxes and the finitely many replacements first.  Uniform
parametric convergence to a simple smooth arc implies that the attachment
points of a replacement of vanishing diameter lie in one vanishing
parameter interval.  The transition points can be chosen in their prescribed
positive-width buffers outside all these intervals.  No bound on the length
or the number of vertices of a replacement is needed.

\subsection{Preliminary separators and genuine stopping explorations}

\begin{lemma}[Many preliminary pairs]\label{lem:preliminary-circuits}
Under the hypotheses of \cref{thm:boundary-arches}, with probability at least
$1-C_0e^{-c_0N}$, at least $c_0N$ of the selected scales have both an open
site circuit at threshold $k-1$ in
\begin{equation}\label{eq:preliminary-strips}
 Q_{\ell_i}\setminus Q_{\ell_i/2}
       \quad\hbox{and one in}\quad
 Q_{7\ell_i}\setminus Q_{5\ell_i},
\end{equation}
each surrounding its inner box.  Artificially open sites are as in
\eqref{eq:boundary-open-sites}.
\end{lemma}

\begin{proof}
Apply the anchored comparison with the high interval on the two clean
portions of the base arc running from just inside the innermost comparison
strip to beyond the outermost strip.  If a central exception is present,
the two portions are separated by a low central portion.  The other
skeleton sites receive the low interval.  The transition buffers are
disjoint from all the strips and their surrounding buffers.  The bounds in
\eqref{eq:arch-boundary-data} give \eqref{eq:synthetic-height-domination}.

Before imposing the complementary transition labels, every pin in a buffer
used for a $t=s$ double circuit is favorable.  By
\cref{lem:buffer-circuits}, there is a constant $a>0$ such that, conditional
on all first-color spins outside a block buffer, a pair of such double
circuits in slightly narrower versions of \eqref{eq:preliminary-strips}
has probability at least $a$.  Choose $M$ so that the block buffers are
disjoint.  Successive revelation of these buffers gives the same lower
bound conditional on the preceding outcomes.  Indeed, for $\lambda>0$ and
the success indicators $I_i$,
\[
 \EE\exp\!\left(-\lambda\sum_{i=1}^N I_i\right)
       \le (1-a+ae^{-\lambda})^N.
\]
It follows that $\sum I_i\ge aN/2$ except with probability $C e^{-cN}$.
Requiring at most four complementary labels multiplies this failure bound
by at most $2^4$.

Each double circuit meets the high part of the skeleton: the original
boundary has a segment from the inner box to the exterior of the outer
box.  The adjacent faces along the double circuit all have the same first
color and, by continuation from that intersection, their lifts lie in the
anchored pair $[k-1,k]$.  The corresponding closed primal walk stays within
a bounded number of lattice spacings of the dual circuit and has winding
number one about the inner box.  It contains a simple surrounding site
circuit.  This produces the preliminary circuits in the synthetic height.
The event that at least a given number of pairs exist is increasing in the
height, so \eqref{eq:synthetic-height-domination} transfers its probability
to the actual law.
\end{proof}

We now select the preliminary pairs by exploring the actual height law.
The synthetic configuration used to bound their abundance plays no part
in this selection.

\begin{lemma}[Extremal-circuit stopping]\label{lem:extremal-stopping}
In a deterministic annular strip, good surrounding simple site circuits
have inner and outer envelopes.  The innermost such circuit can be selected
by revealing the strip from its inner side; the outermost can be selected
from its outer side.  In each case the selection and the revealed values
are measurable with respect to the sites on the circuit and on the tested
side.  No height on the untested side is conditioned by the selection.

In particular, in each successful block of
\cref{lem:preliminary-circuits}, select an innermost circuit $C_i^-$ in the
inner strip and an outermost circuit $C_i^+$ in the outer strip.  Conditional
on these selections and the full configuration outside the intervening free
regions, those regions have independent Gibbs interiors with their
revealed boundary values.
\end{lemma}

\begin{proof}
If two surrounding simple cycles have at least two common vertices, their
union is two-vertex-connected.  After deleting any one vertex, the
remaining part of each cycle is connected and the two parts still share a
vertex.  Its complementary face containing the inner box and its unbounded
face therefore have simple cycle boundaries.  These boundaries are the
inner and outer envelopes.  They use only good vertices of the original
cycles and stay in the same strip.  If there is only one common vertex,
the two cycles cannot cross just once.  Since both enclose the inner box,
their interiors are nested.  Disjoint surrounding cycles are nested as
well.  Iterating the envelope operation in the finite strip gives the
extremal circuit, regarded without an orientation.

For a specified good circuit $C$, whether it is the inner extremal circuit
can be decided from its closed inner side in the strip.  A competing
circuit that improved it while partly leaving that side would have an
inner envelope with $C$ that is itself a better good circuit entirely on
that side.  Thus the absence of a better circuit has no dependence on
unseen heights outside $C$.  The outer statement is identical with the
directions reversed.  Equivalently, these statements specify stopping
sigma-fields: on the event that the selected circuit is $C$, the recorded
heights and the selection event are measurable using the deterministic set
consisting of $C$ and its tested side.  One may reveal that entire set,
including bounded pockets.  Failure to find a circuit requires only the
strip itself.

For the selected pairs, let $V_i$ be the actual free sites between $C_i^-$
and $C_i^+$.  Reveal the unsuccessful strips, the tested sides of successful
strips, the selected circuits, and then every site outside $\bigcup_i V_i$.
Call the resulting sigma-field $\mathcal G$.  On every realization of the
selected sets, the preceding measurability statement says that the event
specifying these sets places no condition on the heights in $\bigcup_i V_i$
beyond their recorded boundary values.  The finite Gibbs specification
therefore applies.  Distinct $V_i$ have no interacting edge after their
separators are fixed, giving conditional independence.  Touches with the
original boundary may divide $V_i$ into several bounded faces; the same
argument applies to their union.
\end{proof}

The selected circuits have height at least $k-1$ on their usable sites.
Every usable boundary value of a retained middle problem is consequently at
least $k-2$.  The full strip
\begin{equation}\label{eq:untouched-middle}
                     Q_{5\ell_i}\setminus Q_{\ell_i}
\end{equation}
between them was not searched.  Its original usable boundary, including any
tiny replacements, still has the clean lower bound $k$.  Moreover, each
selected circuit meets the base polygon boundary, so
$S=S_0\cup C_i^-\cup C_i^+$ is connected.  These facts hold on every
realization of $\mathcal G$ and will be the only boundary information used
to upgrade the block.

\subsection{An attachment from the prescribed boundary}

\begin{lemma}[Boundary attachment]\label{lem:boundary-attachment}
Consider a retained middle problem at scale $\ell$.  Near either of the
locations $(-2\ell,0)$ and $(2\ell,0)$, there are fixed constants
$0<d/W_0\ll1$ and $0<W_0/\ell\ll1$ for which the following holds.
If the local boundary and the tiny replacements are within a sufficiently
small fraction of $d$ of the horizontal boundary line, then with probability
at least $a_0>0$ a path of sites at height at least $k$ enters from the clean
boundary band and reaches depth $d$ inside the box of half-width $W_0$.
Its boundary end is connected through good fixed or artificial sites to
the exterior below, within a slight enlargement of the box.  The bound is
uniform in the exposed boundary values of the middle problem.
\end{lemma}

\begin{proof}
Translate the attachment location to the origin.  Take $W_0$ to be a small
fixed fraction of $\ell$, and later take $d$ to be a smaller fixed fraction
of $W_0$.  The two preliminary circuits are outside this box.  The
unobstructed part above the boundary tolerance belongs to one bounded face
of $S$; in particular, a site $x$ close to $(0,3d)$ is free in that face.

\emph{Synthetic boundary and its mean.}
On $S$, prescribe the high interval $[k-1,k]$ on a base subarc near the
origin, of half-width approximately $W_0/4$, and the low interval
$[k-3,k-2]$ elsewhere.  Choose its endpoints so that the high part covers
horizontal width $W_0/8$, lies within horizontal half-width $W_0/3$, and
avoids all tiny replacement attachments.  A replacement attached within
the high portion receives the high prescription throughout.  There are
only two transition labels.  This construction is feasible for all
sufficiently small replacement diameters and gives the synthetic lower
comparison $K$ in \eqref{eq:synthetic-height-domination}.

Let $A$ be the event that a primal path of heights at least $k$ joins this
high patch to depth $d$ inside the attachment box.  This is an increasing
height event, so it is enough to give it positive probability in the
synthetic law.  We shall first use the mean at $x$ to detect which side of
the boundary-entering contours contains $x$, and then show that the high
side forces $A$.

Between radii of order $d$ and a small fixed fraction of $W_0$, annular
buffers centered on the projected boundary contain only favorable $t=s$
pins.  The annular stack estimate gives, except with probability
\begin{equation}\label{eq:attachment-cut-error}
                      e(d/W_0)\longrightarrow0,
\end{equation}
a double $t=s$ circuit surrounding $x$ and meeting the high patch.  The
same bound, multiplied by a fixed factor, holds after the two complementary
labels are required.  Choose such a cut using only the first color.  Its
inner face chain belongs to the anchored pair $[k-1,k]$, and both transition
labels lie outside it.  No component of $t$-changing edges crosses the
double cut.  Thus the complementary labels in its interior, including its
inner cut layer, are all free.  Flipping these labels preserves their
conditional law.  It reverses every interior height increment and swaps
the two values of the anchored cut pair, hence sends
$K$ to $2k-1-K$ there.  The conditional mean at $x$ is therefore $k-1/2$.

We also need a bound on the mean on failure, uniform in the lattice size.
The component potentials in the proof of \cref{lem:component-switch}
give this bound.  On each component of $t$-changing edges, $dt$ extended
by zero has a potential with two values differing by two.  If $t=B$, the
identity
\[
                 d(K+BW/2)=W\,dB
\]
expresses the height difference from the transition anchor to $x$ as an
endpoint term of absolute value at most one, plus the component labels
times their potential differences.  Each potential difference has absolute
value at most two.  If $t=W$, use $d(K-BW/2)=-B\,dW$ instead.
The unprescribed labels have conditional mean zero, and at most two
components have prescribed labels.  Since the anchor has value $k-2$ or
$k-1$, this proves
\begin{equation}\label{eq:attachment-failure-mean}
       \bigl|\EE[K(x)\mid t,\text{transition labels}]-k\bigr|\le C_1
\end{equation}
for every first-color configuration, with an absolute constant $C_1$.
Combining the symmetry on success with this bound gives
\begin{equation}\label{eq:attachment-synthetic-mean}
                  \EE K(x)=k-\tfrac12+O(e(d/W_0)).
\end{equation}
All these arguments take place in a finite free volume around the
construction, so no infinite-volume height or absolute anchor is required.

\emph{Low lateral barriers.}
Around the boundary locations with horizontal offsets $\pm2W_0/3$, use
annuli with inner radius several times $d$ and outer radius less than
$W_0/8$.  All local first-color pins are now favorable to $t=-s$.  With
failure probability again tending to zero with $d/W_0$, each side contains
a double $t=-s$ circuit meeting the low skeleton.  Its face paths have the
anchored low values and include a path from the low skeleton to above
depth $3d$, inside the horizontal band of width $W_0/4$ about that side
location.  Write $L$ for the event that both such lateral barriers exist.
They have height at most $k-2$ in the synthetic problem.
The estimates survive the same two complementary labels.  We may increase
the lower end of the annular range by a fixed factor when a little more
vertical clearance is needed.

For the planar separation below, continue each barrier along its local low
skeleton arc and then through the exterior to a line below the boundary
band.  Uniform parametric approximation keeps this continuation in a
slightly enlarged lateral strip; a tiny replacement attaches to the same
local base arc.  The skeleton portion is outside the high patch and is
low throughout.  The exterior portion is only a geometric continuation,
which a contour lying in the face cannot cross.  Thus the two barriers
separate the high patch laterally all the way through the boundary band.

\emph{The mean detects a side of the open contours.}
Condition on the synthetic trace on the whole skeleton $S$.  In its
bounded face containing $x$, follow every dual level path at threshold
$k-1/2$ that enters through a skeleton edge.  At each crossed edge, the
adjacent heights are exactly $k-1,k$.  On entering a triangle, those two
values force its third value to be one of $k-1,k$, and hence determine the
next crossed edge.  The exploration only reveals the adjacent values of
the paths being followed.  Exhaust all boundary crossings in this way,
including crossings on either side of a skeleton edge bordering a slit
in the face.

Conditional on these open paths and their lips, their occurrence imposes
only the indicated fixed vertex values.  Any admissible completion with
those values has the same open paths, by the degree-two rule and exhaustion
of the boundary crossings.  Unseen closed contours remain unrestricted.
The remaining free components therefore have their Gibbs laws.

Cut the face along the exposed open paths.  Their lips have known sides:
the $k$ lip is high and the $k-1$ lip is low.  Together with the threshold
signs of the skeleton trace, these sides label every resulting geometric
cell high or low.  Here is why these labels are consistent even when the
skeleton boundary has slits.  A closed threshold contour in this face
cannot enclose a skeleton site, since the connected skeleton continues
outside the contour and cannot cross it.  It cannot enclose part of an
exposed open path either: that path continues to the skeleton and contours
do not intersect.  In the full threshold-sign picture, one may therefore
erase the closed contours, filling each outermost enclosed region with
its exterior sign, without changing any recorded skeleton or lip sign.
The remaining sign is constant in each cell cut out by the open paths.
This proves consistency; it is a deterministic explanation, not an
additional revelation or conditioning of the closed contours.

The free primal Gibbs components left by the exploration lie in these
geometric cells; removing the revealed lip vertices may divide a cell
further.  In a low cell every adjacent fixed value is at most $k-1$.
In a high cell every adjacent fixed value is at most $k$: the high
skeleton values and the high lips are exactly $k$, and all other fixed
values satisfy the same ceiling.  Let $U$ be the event that $x$ belongs
to a high cell, including the case that it lies on a revealed high lip.
A revealed low lip belongs to the low case.

Comparison with a constant ceiling, including the deterministic case when
$x$ lies on a revealed lip, yields
\begin{equation}\label{eq:attachment-cell-mean}
 \EE K(x)\le k-1+\PP(U).
\end{equation}
Together with \eqref{eq:attachment-synthetic-mean}, this gives
$\PP(U)\ge1/2-O(e(d/W_0))$.  A closed high island around $x$ does not
suffice for $U$.  Such islands remain inside their unrestricted Gibbs
cells, and comparison controls the cell's mean rather than each of its
height values.

\emph{The correct boundary lip reaches the box.}
The high lip of every exposed open path is a connected primal chain at
height $k$.  Its endpoints are skeleton sites at height $k$, so they lie
on the designated high patch, including any tiny replacements carrying
that patch's prescription.  No endpoint lies on a preliminary circuit or
on another low part of the skeleton.

We claim that $U\cap L\subseteq A$.  Suppose instead that $L$ occurs but
$A$ does not.  Follow a high lip from either of its endpoints on the high
patch.  It cannot cross a low lateral barrier: these are primal paths of
height at most $k-2$, whereas the lip has height $k$.  It cannot leave
through the exterior side of the face.  To pass above a barrier it would
first have to reach depth $d$, already giving $A$.  Consequently every
exposed open path stays between the two barriers and below depth $d$.
This statement concerns all the open paths, since every one of their high
lips starts on the designated high patch.

The site $x$ is at depth $3d$.  Staying above depth $2d$, it can be
joined to one lateral barrier without crossing any exposed open contour.
Following that barrier down to its low skeleton endpoint also crosses
none: all its sites have height at most $k-2$.  The resulting route puts
$x$ in the cell with the low label inherited from that endpoint.  This
contradicts $U$.  The boundary and the tiny replacements stay in a much
shallower band and do not obstruct the deep portion of this route; the
barrier itself supplies the continuation through that band.  Fixed margins
between depths $d$ and $3d$ absorb lattice rounding.

Since $\PP(L^c)=O(e(d/W_0))$, the inclusion proves
\[
 \PP(A)\ge\PP(U)-\PP(L^c)\ge\tfrac12-O(e(d/W_0)).
\]
The mean calculation was not conditioned on the lateral barriers.
Choose $d/W_0$ once so that this bound is at least, say, $1/4$.
The event $A$ asks only for a superlevel path, even though the synthetic
construction produced an exact-height-$k$ lip.  It is therefore increasing,
and \eqref{eq:synthetic-height-domination} gives the same positive lower
bound for the actual height.

Finally, in the actual comparison every usable skeleton site in the local
box is at least $k$; the other skeleton sites there are artificial.  A tiny
replacement is attached to the local base arc, and that arc, with its
parametric slack at the box margins, is connected through open sites to
the exterior below.  For example a path from below stopped at its first
encounter with the base boundary reaches such an open site.  Thus the
attachment has the exterior continuation claimed in the lemma.
\end{proof}

\subsection{Joining the attachments and amplifying across scales}

\begin{lemma}[Upgrade in a retained middle region]\label{lem:arch-upgrade}
For every successful preliminary block, conditional on $\mathcal G$ from
\cref{lem:extremal-stopping}, the probability of an arch separator at
threshold $k$ in its untouched middle is at least an absolute constant
$a_1>0$.
\end{lemma}

\begin{proof}
Fix the realized middle problem.  Its actual boundary has the lower bounds
recorded after \eqref{eq:untouched-middle}.  The two attachments of
\cref{lem:boundary-attachment} have uniform positive probabilities, and
their events are increasing in the actual height.

We first record a bulk path estimate under this same conditional law.
For a box with a fixed relative annular buffer inside the free middle,
compare from below with constant height $k-2$ on its bounding skeleton.
The latter comparison may be realized by fixing both spins of that height
on $S$.  Connectedness of $S$ fixes a common flat lift on every relevant
bounded face.  Its height contours are closed loops: a contour cannot
cross the flat skeleton or enclose just part of a connected skeleton.
Given the unoriented loops, all their jump signs are independent fair signs.
For use of the first-spin circuit estimates, first enforce constancy of
the second spin along $S$.  This is an $E_0$ constraint outside the bulk
buffers.  Since $S$ is connected, specifying its second-spin value costs
one fair component label and does not change the first-spin marginal.

Apply \cref{lem:buffer-circuits} in a fixed number of thinner concentric
buffers, requesting alternating first-spin signs on five double circuits.
All these buffers are inside the free region.  Successive conditional
estimates give a positive probability, uniform in mesh and exposed actual
heights, for this event in the flat comparison.  Between each successive
pair there is a first-spin wall loop surrounding the inner box.  Hence
there are at least four nested height loops in the allocated annulus.

Retain only the event that there are at least four enclosing unoriented
loops in the allocated annulus, forgetting the first-spin event used to
lower-bound its probability.  This event and the choice of four consecutive
enclosing loops depend only on the full unoriented geometry.  Conditional
on that geometry, the height just outside the first chosen loop is $k-2$
plus a sum of other fair loop signs.  It is therefore symmetric about
$k-2$ and independent of the four chosen signs.  With probability at least
$2^{-5}$ it is at least $k-2$ and all four signs are positive.  The inner
face path of the fourth loop is then at least $k$.  Its closed primal walk
has winding number one and contains a surrounding circuit, with only
bounded mesh displacement.  Thus a fixed bulk annulus has a uniformly
positive probability of a circuit at height at least $k$.  Height order
transfers this to the actual middle law.

These bulk circuits give paths along any fixed polygonal tube with positive
width and clearance.  To see the deterministic gluing, place small annuli
along the tube so that consecutive inner disks overlap and each protrudes
beyond the other annulus's outer disk.  Two surrounding circuits cannot
then be disjoint or nested, so they intersect.  A finite string supplies
the desired connected crossing.  Positive association for heights bounds
the probability of the simultaneous increasing circuit events below by the
product of their positive lower bounds.

Use such a tube to cross each attachment box horizontally near depth
$d/2$, extending a little beyond its sides.  Any rising attachment from
the boundary band to depth $d$ intersects this crossbar.  Join the left
and right crossbars by tubes going upward near horizontal locations
$\pm2\ell$, across at depth $2\ell$, and downward on the right.  Tube
widths are much smaller than $d$, and all their bulk boxes have clearance
from the skeleton.  The whole construction lies between $Q_\ell$ and
$\partial Q_{5\ell}$; see \cref{fig:boundary-arches}.

Apply height positive association in the actual conditional Gibbs problem
to the two attachment events and the finite collection of bulk circuit
events.  This gives a fixed lower bound $a_1>0$ for their intersection.
The exterior continuations of the attachments can be joined below the
boundary, around the bottom of $Q_\ell$, through artificial sites and
good local boundary sites.  The resulting closed open walk follows a
rectangle around $Q_\ell$ and has winding number one, so it contains an
open surrounding simple cycle.  The tubes avoid the preliminary circuits;
in particular, the construction has not declared any of their possibly
sub-$k$ usable values artificial.  Every artificial site used here was
already outside the original region or in a removed boundary patch,
which proves the upgrade.
\end{proof}

\begin{figure}[tb]
 \centering
 \begin{tikzpicture}[x=0.42cm,y=0.42cm,>=Stealth,font=\small]
  \fill[gray!12,even odd rule] (-7,-7) rectangle (7,7)
                              (-5,-5) rectangle (5,5);
  \fill[gray!24,even odd rule] (-1,-1) rectangle (1,1)
                              (-.5,-.5) rectangle (.5,.5);
  \fill[blue!3] (-5,0) rectangle (5,5);
  \fill[gray!24,even odd rule] (-1,-1) rectangle (1,1)
                              (-.5,-.5) rectangle (.5,.5);
  \draw[gray!60] (-7,-7) rectangle (7,7);
  \draw[gray!60] (-5,-5) rectangle (5,5);
  \draw[gray!65] (-1,-1) rectangle (1,1);
  \draw[gray!65] (-.5,-.5) rectangle (.5,.5);
  \node[anchor=east,font=\scriptsize] at (-7.35,7) {$Q_{7\ell}$};
  \draw[gray!65,thin] (-7.25,7)--(-7,7);
  \node[anchor=east,font=\scriptsize] at (-7.35,5) {$Q_{5\ell}$};
  \draw[gray!65,thin] (-7.25,5)--(-5,5);
  \node[anchor=center,font=\scriptsize,fill=white,inner sep=.5pt] at (0,1.5) {$Q_\ell$};
  \draw[gray!65,thin] (0,1.1)--(0,1);
  \node[anchor=center,font=\scriptsize] at (0,-1.45) {$Q_{\ell/2}$};
  \draw[gray!65,thin] (0,-1.12)--(0,-.5);
  \draw[black!65,thick] (-7.35,0)--(7.35,0);
  \node[anchor=south east,font=\scriptsize] at (7.25,.15) {clean boundary};
  \node[blue!65!black] at (0,4) {free middle};
  \node[gray!70!black] at (0,-4) {artificial exterior};
  \node[font=\scriptsize] at (0,6.05) {outer strip};
  \node[anchor=east,font=\scriptsize,fill=white,inner sep=1pt] at (-2.7,-1.1) {inner strip};
  \draw[gray!65,thin] (-2.65,-1.1)--(-1,-.8);
  \draw[blue!70!black,line width=1.3pt] (-2,.08)--(-2,2)--(2,2)--(2,.08);
  \draw[blue!70!black,line width=1.3pt] (-2.35,.35)--(-1.65,.35);
  \draw[blue!70!black,line width=1.3pt] (1.65,.35)--(2.35,.35);
  \draw[blue!70!black,dashed,line width=1.1pt]
       (-2,.02)--(-2,-2)--(2,-2)--(2,.02);
  \draw[orange!80!black,line width=1.2pt] (-2.1,.02)--(-1.95,.65);
  \draw[orange!80!black,line width=1.2pt] (1.9,.02)--(2.05,.65);
  \draw[black!65,densely dotted] (-2.5,-.25) rectangle (-1.5,.85);
  \draw[black!65,densely dotted] (1.5,-.25) rectangle (2.5,.85);
  \fill (0,0) circle (.065);
  \node[anchor=north west,font=\scriptsize] at (.08,-.02) {$p$};
  \node[anchor=north,font=\scriptsize,fill=white,inner sep=1pt] at (-2,-.3) {$q_L$};
  \node[anchor=north,font=\scriptsize,fill=white,inner sep=1pt] at (2,-.3) {$q_R$};
  \node[anchor=south,font=\scriptsize] at (0,2.12) {bulk joining tube};
  \begin{scope}[shift={(12,-.15)},x=1.4cm,y=0.65cm]
    \fill[blue!3] (-1.1,0) rectangle (1.1,3.65);
    \draw[black!60] (-1.15,0)--(1.15,0);
    \draw[orange!80!black,line width=2pt] (-.3,0)--(.3,0);
    \node[anchor=north,font=\scriptsize] at (0,-.13) {anchored high patch};
    \draw[red!70!black,line width=1.1pt] (-.70,0)--(-.73,1.1)--(-.64,2.1)--(-.68,3.45);
    \draw[red!70!black,line width=1.1pt] (.70,0)--(.65,1.2)--(.74,2.2)--(.68,3.45);
    \node[rotate=90,font=\scriptsize,red!65!black] at (-.83,1.85) {low barrier};
    \node[rotate=90,font=\scriptsize,red!65!black] at (.88,1.85) {low barrier};
    \draw[orange!80!black,line width=1.3pt] (-.1,0).. controls (.2,.3) and (-.18,.8)..(.12,1.15);
    \draw[gray!70,densely dashed] (-1.05,1)--(1.05,1);
    \node[anchor=west,font=\scriptsize] at (1.09,1) {$d$};
    \fill (0,3) circle (.035);
    \node[anchor=north,align=center,font=\scriptsize] at (0,2.8) {$x$ at depth\\$3d$};
    \draw[gray!60,densely dotted] (-1.04,-.05) rectangle (1.04,3.6);
    \node[align=center,font=\scriptsize] at (0,4.3) {Synthetic attachment\\vertical scale magnified};
  \end{scope}
\end{tikzpicture}


 \caption{A retained block in tangent coordinates.  The four labeled squares
 determine the two searched strips; the attachment boxes and path widths are
 enlarged for legibility.  The two shaded strips contain the preliminary separators.
 Their facing sides leave the middle Gibbs region unsearched.  The blue
 route uses two boundary attachments, horizontal crossbars, and bulk
 paths at positive clearance in the actual height law.  The right detail
 shows the synthetic low-barrier argument used to prove an attachment;
 its barriers are not imposed in the actual law.  The dashed lower route
 uses the artificial exterior.  All displayed paths are schematic corridors,
 not a prescribed realization of the random separator.}
 \label{fig:boundary-arches}
\end{figure}

\begin{proof}[Proof of \cref{thm:boundary-arches}]
Let $Z$ be the number of successful preliminary blocks.
\Cref{lem:preliminary-circuits} gives
$\PP(Z<c_0N)\le C_0e^{-c_0N}$.  Use the actual-height stopping sigma-field
$\mathcal G$ of \cref{lem:extremal-stopping}.  Conditional on it, the retained
middle regions have independent Gibbs laws, and
\cref{lem:arch-upgrade} gives success probability at least $a_1$ in each.
Consequently
\[
 \PP(\text{no upgraded block})
 \le C_0e^{-c_0N}+(1-a_1)^{c_0N}
 \le Ce^{-cN}.
\]
Every upgraded circuit surrounds $Q_r$ and lies in $Q_R$ by
\eqref{eq:arch-scales}.  This proves
\eqref{eq:arch-high-probability}.  Applying the argument to $-H$ proves the
reversed assertion.
\end{proof}

\begin{corollary}[Random traces and tiny boundary replacements]
\label{cor:boundary-robustness}
The arch estimate remains valid conditional on an exterior sigma-field
whenever, conditional on that field, the remaining free heights have a
uniform Gibbs law and their actual usable boundary data obey
\eqref{eq:arch-boundary-data}.  It also remains valid after a fixed finite
number of boundary-attached patches have been removed by exposed simple
circuits, provided the same data bounds hold on the new usable boundary.

For a prescribed error, first fix a finite collection of arch scales and
their boxes.  The replacement diameters, boundary approximation errors, and
mesh may then be made sufficiently small for all comparisons.  The success
constants of the resulting blocks have no factor depending on the number
of microscopic pins or on the fixed number of replacements.
\end{corollary}

\begin{proof}
Every probabilistic lower bound in the proof is uniform in the specified
trace; only its displayed one-sided bounds were used.  It may therefore be
applied to each conditional trace.  A boundary-attached replacement keeps
the skeleton connected and preserves the full-neighborhood condition for
the retained free sites.  For fixed scales, sufficiently small replacements
stay in the shallow boundary tolerance and do not meet the bulk tubes.
The transition points of the synthetic prescriptions can be chosen away
from all attachment neighborhoods, as explained after
\eqref{eq:synthetic-height-domination}.  The number of complementary labels
is consequently still at most four in the preliminary construction and at
most two in each attachment construction.  These are the only label costs
in the proof.  The geometric approximation threshold may depend on the
finite family, but none of the block probability constants does.
\end{proof}
